\documentclass[12pt,a4paper]{article}

\usepackage[utf8]{inputenc}
\usepackage[T1]{fontenc}
\usepackage[french]{babel}
\usepackage{geometry}
\usepackage{graphicx}
\usepackage{longtable}
\usepackage{array}
\usepackage{booktabs}
\usepackage{amsmath}
\usepackage{hyperref}
\usepackage{float}
\usepackage{listings}
\usepackage{xcolor}
\usepackage{titlesec}
\usepackage{enumitem}

\geometry{margin=2.5cm}

\hypersetup{
    colorlinks=true,
    linkcolor=blue,
    urlcolor=blue
}

\title{
\vspace{-2cm}
\textbf{\Huge Système de Détection de Spam et Phishing}\\
\vspace{0.3cm}
\Large Basé sur l'IA et les Indicateurs de Compromission (IoCs)
}

\author{
\textbf{Projet Cybersécurité - Personne A}\\
Analyse des Menaces et Ingénierie des IoCs\\
Encadré par Pr. Sara Ouald Chaib
}

\date{2026}

\begin{document}

\maketitle

\tableofcontents
\newpage

% =========================================================
\section{Introduction}

L'objectif principal de ce projet est de concevoir un système intelligent capable de détecter les emails malveillants à travers une approche orientée cybersécurité.

Contrairement à une approche purement centrée sur le Machine Learning, ce projet considère le spam et le phishing comme de véritables cyberattaques exploitant des mécanismes techniques et psychologiques afin de compromettre un système d'information.

Le système proposé vise donc à :
\begin{itemize}
    \item Identifier les emails frauduleux
    \item Détecter les tentatives de phishing
    \item Analyser les comportements malveillants
    \item Construire des Indicateurs de Compromission (IoCs)
    \item Fournir des données pertinentes aux modèles de détection
\end{itemize}

Cette approche s'aligne directement avec les concepts étudiés dans le module de cybersécurité :
\begin{itemize}
    \item Ingénierie sociale
    \item Phishing
    \item Filtrage réseau
    \item Architecture défensive
    \item Détection d'intrusions
    \item Analyse comportementale
\end{itemize}

% =========================================================
\section{Analyse des Menaces}

\subsection{Cybermenace vs Cyberattaque}

Une cybermenace représente un danger potentiel permanent pouvant affecter un système d'information.

Une cyberattaque représente le passage à l'action. Elle vise généralement :
\begin{itemize}
    \item La confidentialité
    \item L'intégrité
    \item La disponibilité
\end{itemize}

Dans le contexte des emails malveillants, les attaques exploitent principalement :
\begin{itemize}
    \item La manipulation psychologique
    \item L'usurpation d'identité
    \item Les redirections malveillantes
    \item Les pièces jointes piégées
\end{itemize}

% =========================================================
\subsection{Spam Traditionnel}

Le spam correspond à l'envoi massif de messages non sollicités.

\subsubsection{Objectifs du Spam}

Les objectifs principaux sont :
\begin{itemize}
    \item Saturation des systèmes
    \item Publicité frauduleuse
    \item Diffusion massive de liens
    \item Escroqueries simples
\end{itemize}

\subsubsection{Impact Sécurité}

Le spam impacte principalement :
\begin{itemize}
    \item La disponibilité du système
    \item Les performances réseau
    \item La productivité des utilisateurs
\end{itemize}

Le risque de compromission directe reste relativement faible comparé au phishing.

% =========================================================
\subsection{Phishing et Spear Phishing}

Le phishing est une technique d'ingénierie sociale visant à tromper la victime afin d'obtenir :
\begin{itemize}
    \item Des identifiants
    \item Des données bancaires
    \item Des informations sensibles
    \item Une exécution de malware
\end{itemize}

\subsubsection{Ingénierie Sociale}

L'ingénierie sociale exploite :
\begin{itemize}
    \item La peur
    \item L'urgence
    \item La confiance
    \item L'autorité
\end{itemize}

L'objectif est de pousser l'utilisateur à agir rapidement sans vérifier l'authenticité du message.

\subsubsection{Spear Phishing}

Le spear phishing est une version ciblée du phishing.

L'attaquant personnalise l'email pour :
\begin{itemize}
    \item Un employé précis
    \item Un directeur
    \item Une organisation spécifique
\end{itemize}

Ces attaques sont beaucoup plus difficiles à détecter.

% =========================================================
\section{Comportement Adaptatif des Attaquants}

Les cybercriminels adaptent constamment leurs techniques afin de contourner les systèmes de sécurité.

Les filtres statiques basés uniquement sur des mots-clés deviennent rapidement inefficaces.

Les attaquants utilisent notamment :
\begin{itemize}
    \item L'obfuscation Unicode
    \item Les caractères invisibles
    \item Les domaines générés aléatoirement
    \item Les homoglyphes
    \item Les redirections multiples
    \item Le spoofing
\end{itemize}

Le système doit donc détecter :
\begin{itemize}
    \item Les anomalies structurelles
    \item Les comportements suspects
    \item Les indicateurs comportementaux
\end{itemize}

et non uniquement des signatures fixes.

% =========================================================
\section{Architecture de Défense Multi-Couches}

Le système d'analyse des emails a été conçu comme une architecture de filtrage multi-couches.

Chaque couche joue un rôle spécifique dans la détection des menaces.

\subsection{Principe du Filtrage Multi-Couches}

Le principe consiste à analyser simultanément plusieurs surfaces d'attaque :

\begin{enumerate}
    \item Les headers
    \item Les URLs
    \item Le contenu textuel
    \item Les comportements psychologiques
    \item Les métadonnées
    \item Les pièces jointes
\end{enumerate}

Même si une couche est contournée, les autres couches peuvent toujours détecter l'attaque.

% =========================================================
\section{Indicateurs de Compromission (IoCs)}

Les Indicateurs de Compromission représentent des éléments observables permettant d'identifier une activité potentiellement malveillante.

Dans ce projet, les IoCs sont utilisés comme caractéristiques d'entrée pour les modèles de Machine Learning.

% =========================================================
\subsection{Header-Based IoCs}

Ces indicateurs analysent les anomalies liées au routage et à l'authentification des emails.

\subsubsection{Reply-To Mismatch}

Le Reply-To Mismatch correspond à une différence entre :
\begin{itemize}
    \item L'adresse d'expédition
    \item L'adresse de réponse
\end{itemize}

Cette technique est fréquemment utilisée dans les attaques de spoofing.

% =========================================================
\subsubsection{SPF, DKIM et DMARC}

\paragraph{SPF}
Le SPF (Sender Policy Framework) permet de vérifier si le serveur expéditeur est autorisé à envoyer des emails pour un domaine donné.

\paragraph{DKIM}
Le DKIM (DomainKeys Identified Mail) ajoute une signature cryptographique permettant de vérifier l'intégrité du message.

\paragraph{DMARC}
Le DMARC définit la politique appliquée lorsqu'un email échoue aux vérifications SPF ou DKIM.

Un défaut d'alignement SPF/DKIM/DMARC constitue un IoC critique.

% =========================================================
\subsection{URL-Based IoCs}

Les URLs représentent l'un des principaux vecteurs d'attaque dans les campagnes de phishing.

% =========================================================
\subsubsection{IP-Based URLs}

Les attaquants utilisent parfois directement des adresses IP afin de contourner :
\begin{itemize}
    \item Les filtres DNS
    \item Les listes noires
    \item Les systèmes de réputation
\end{itemize}

Exemple :
\begin{verbatim}
http://192.168.1.25/login
\end{verbatim}

% =========================================================
\subsubsection{TLDs Suspects}

Certains domaines génériques sont fréquemment utilisés dans les campagnes malveillantes :
\begin{itemize}
    \item .xyz
    \item .top
    \item .click
    \item .shop
\end{itemize}

Ces domaines constituent donc des indicateurs de risque.

% =========================================================
\subsubsection{Homoglyphes}

Les homoglyphes consistent à remplacer certains caractères par des symboles visuellement similaires.

Exemple :
\begin{verbatim}
paypaI.com
\end{verbatim}

Le caractère "I" majuscule remplace la lettre "l".

Cette technique vise à tromper visuellement la victime.

% =========================================================
\subsubsection{Entropie des Domaines}

Les domaines générés automatiquement possèdent souvent une forte entropie.

Exemple :
\begin{verbatim}
xj82ks91-top-security.xyz
\end{verbatim}

Une forte entropie peut indiquer :
\begin{itemize}
    \item Un DGA (Domain Generation Algorithm)
    \item Une infrastructure malveillante
\end{itemize}

% =========================================================
\section{Behavioral and Psychological IoCs}

Ces indicateurs visent à détecter les mécanismes d'ingénierie sociale.

% =========================================================
\subsection{Densité d'Urgence}

Les emails malveillants utilisent fréquemment :
\begin{itemize}
    \item urgent
    \item immédiatement
    \item action requise
    \item compte suspendu
\end{itemize}

L'objectif est de provoquer :
\begin{itemize}
    \item Stress
    \item Panique
    \item Réaction impulsive
\end{itemize}

% =========================================================
\subsection{Credential Requests}

Les emails de phishing cherchent souvent à récupérer :
\begin{itemize}
    \item login
    \item password
    \item OTP
    \item données bancaires
\end{itemize}

La présence répétée de ce vocabulaire représente un indicateur important.

% =========================================================
\section{Content and Obfuscation IoCs}

% =========================================================
\subsection{Ratio HTML/Texte}

Les emails malveillants utilisent souvent :
\begin{itemize}
    \item Beaucoup de HTML
    \item Des boutons cachés
    \item Des scripts
    \item Des redirections
\end{itemize}

Un ratio HTML anormal peut révéler :
\begin{itemize}
    \item Une tentative de dissimulation
    \item Une redirection cachée
\end{itemize}

% =========================================================
\subsection{Majuscules Excessives}

Les campagnes frauduleuses utilisent fréquemment :
\begin{itemize}
    \item Des majuscules massives
    \item Des caractères spéciaux
    \item Des ponctuations excessives
\end{itemize}

Exemple :
\begin{verbatim}
URGENT!!! VERIFY YOUR ACCOUNT NOW!!!
\end{verbatim}

% =========================================================
\subsection{Obfuscation Unicode}

Les attaquants utilisent :
\begin{itemize}
    \item Zero-width characters
    \item Leetspeak
    \item Unicode invisible
\end{itemize}

afin de contourner :
\begin{itemize}
    \item Les filtres lexicaux
    \item Les signatures statiques
\end{itemize}

% =========================================================
\section{Metadata IoCs}

% =========================================================
\subsection{Extensions Dangereuses}

Certaines extensions sont fortement associées aux malwares :
\begin{itemize}
    \item .exe
    \item .vbs
    \item .scr
    \item .js
\end{itemize}

Ces pièces jointes peuvent agir comme :
\begin{itemize}
    \item Droppers
    \item Trojans
    \item Loaders
    \item Ransomwares
\end{itemize}

% =========================================================
\section{Mapping Attaque $\rightarrow$ IoC $\rightarrow$ Défense}

\begin{table}[H]
\centering
\begin{tabular}{|p{4cm}|p{4cm}|p{4cm}|}
\hline
\textbf{Technique d'Attaque} & \textbf{Vecteur} & \textbf{IoC Défensif} \\
\hline
Spoofing & Faux email & Reply-To mismatch \\
\hline
Credential Theft & Faux login & Credential keywords \\
\hline
Malware Delivery & Pièce jointe & Suspicious extension \\
\hline
DNS Bypass & URL IP directe & IP-based URL \\
\hline
Social Engineering & Urgence psychologique & Urgency density \\
\hline
Filter Evasion & Obfuscation Unicode & Unicode anomaly detection \\
\hline
\end{tabular}
\caption{Mapping Attaque $\rightarrow$ IoC $\rightarrow$ Défense}
\end{table}

% =========================================================
\section{Pipeline de Détection}

Le pipeline complet du système suit les étapes suivantes :

\begin{enumerate}
    \item Réception de l'email
    \item Parsing des headers
    \item Extraction des URLs
    \item Analyse comportementale
    \item Extraction des IoCs
    \item Vectorisation
    \item Classification ML
    \item Décision finale
\end{enumerate}

% =========================================================
\section{Conclusion}

Cette approche permet de transformer un simple système de classification en une véritable architecture défensive orientée cybersécurité.

Le projet ne repose pas uniquement sur des modèles statistiques, mais sur :
\begin{itemize}
    \item Une compréhension des cyberattaques
    \item Une analyse comportementale
    \item Des mécanismes de filtrage multi-couches
    \item Des Indicateurs de Compromission réalistes
\end{itemize}

L'objectif final est de construire un système capable de résister aux comportements adaptatifs des attaquants tout en maintenant un faible taux de faux positifs et de faux négatifs.

\end{document}